\section{Infidels and the Faithful}

According to \textbf{Rene Guenon}, theological, or exoteric, statements can be expressed in esoteric terms. In \emph{Symbolism of the Cross}, he illustrates how the esoteric version deepens our understanding with a precise formulation.

In the example he uses from Islam, the categories of the faithful and the infidel cannot depend simply on the nominal affiliation to an exoteric religion. He points out that ``Islam'' means ``submission to God's Will''. Hence, the faithful, or the elect, are those who ``consciously and voluntarily conform to the universal order''.

The infidels are of two types. The rejected ``obey the law against their will'', that is, they oppose the Divine Plan and face God's wrath. Those who have gone astray --- most people --- are ignorant of God's Will and are absorbed in the affairs of the world.

He relates those tendencies to the three gunas.

\begin{itemize}

\item The \textbf{elect} are sattvic, being pulled upward. They follow the ``straight path'', which is the place of ``peace'', and their conformity to the universal Will makes them true collaborators in the Divine Plan.

\item The \textbf{rejected} are tamasic following the downward force along the vertical axis, towards the lower states. Because they do so consciously, this path is the ``infernal'' way opposed to the ``heavenly'' way.

\item The \textbf{ignorant} are rajasic, interested in working out their potentials horizontally in the human state.
\end{itemize}


Since everything is subject to the three gunas, it is necessary, in esoteric work, to learn to recognize their effects. Sattva is the drive to higher states of being and to knowledge. Saints, Prophets, and Sages provide signposts in this direction.

The tamas leads downward from the human state to the animalistic and even material states. Animal life is dominated by the three lowest chakras: fear, sex, and hunger. These appear as anger, lust, and the desire for power. Hence, anything that excites those feelings can be recognized as tamasic.

The majority of people are dominated by rajas, which is recognized by an excessive concern for worldly affairs. Even these are subject to sattva and tamas, but seldom in a fully conscious way. They may experience a call to spirituality, but it is usually in a sentimental way. Others may follow their lower forces, not to consciously oppose God's will, but rather to improve their lot in life.

\paragraph{Appendix}


To be specific about the sattva pull, but not very specific, here are some general recommendations. Much of this material has been covered, but more is forthcoming.

\begin{itemize}


\item Follow your exoteric teaching

\item Learn the universal order

\item Dominate the subhuman states

\item Purify the human states

\item Develop awareness of higher states
\end{itemize}



\flrightit{Posted on 2021-12-29 by Cologero}

